\documentclass[12pt,a4paper]{ctexart}

% ===== 全国大学生数学建模竞赛（CUMCM）论文模板 =====
% 国赛格式规范：
%   - A4纸，上下左右各2.5cm
%   - 标题：三号黑体加粗居中
%   - 摘要：小四宋体，"摘要"二字小四黑体
%   - 一级标题：四号黑体，数字编号(1,2,3...)
%   - 正文：小四宋体，首行缩进2字符
%   - 表格：三线表
%   - 页码：底部居中，从正文第1页编起

% ---------- 页面设置 ----------
\usepackage[top=2.5cm,bottom=2.5cm,left=2.5cm,right=2.5cm]{geometry}

% ---------- 数学与公式 ----------
\usepackage{amsmath,amssymb,amsthm}

% ---------- 图表 ----------
\usepackage{graphicx,subfig}
\usepackage{booktabs,multirow,array,longtable}
\usepackage{float}

% ---------- 算法伪代码 ----------
\usepackage{algorithm,algorithmic}
\renewcommand{\algorithmicrequire}{\textbf{输入:}}
\renewcommand{\algorithmicensure}{\textbf{输出:}}

% ---------- 引用与链接 ----------
\usepackage[colorlinks=true,linkcolor=black,citecolor=black,urlcolor=black]{hyperref}

% ---------- 列表 ----------
\usepackage{enumitem}

% ---------- 图表标题 ----------
\usepackage{caption}
\captionsetup{font=small,labelfont=bf,labelsep=quad,justification=centering}

% ---------- 页眉页脚 ----------
\usepackage{fancyhdr}
\pagestyle{fancy}
\fancyhf{}
\fancyhead[C]{\songti\zihao{5}\leftmark}     % 页眉：居中显示当前节名
\fancyfoot[C]{\thepage}                        % 页脚：居中页码
\renewcommand{\headrulewidth}{0.4pt}
\setlength{\headheight}{14pt}

% ---------- 首行缩进 ----------
\usepackage{indentfirst}
\setlength{\parindent}{2em}

% ---------- 行距 ----------
\linespread{1.3}

% ---------- 标题格式 ----------
% 国赛使用数字编号：1, 1.1, 1.1.1
\CTEXsetup[name={,},number={\arabic{section}},format={\heiti\zihao{4}\centering},aftername={}]{section}
\CTEXsetup[name={,},number={\arabic{section}.\arabic{subsection}},format={\heiti\zihao{-4}\raggedright},aftername={}]{subsection}
\CTEXsetup[name={,},number={\arabic{section}.\arabic{subsection}.\arabic{subsubsection}},format={\heiti\zihao{-4}\raggedright},aftername={}]{subsubsection}

% 图表编号按节
\numberwithin{equation}{section}
\numberwithin{figure}{section}
\numberwithin{table}{section}
\renewcommand{\figurename}{图}
\renewcommand{\tablename}{表}

\begin{document}

% ==========================================
% 首页：承诺书区域（国赛特有）
% ==========================================
% 国赛论文首页需要包含参赛承诺书，此处仅留提示
% 实际比赛时请使用组委会提供的模板首页

\vspace*{2cm}
\begin{center}
{\heiti\zihao{3}\textbf{【在此处填写论文标题】}}

\vspace{1.5cm}

{\heiti\zihao{4}摘\quad 要}
\end{center}

\vspace{0.5cm}

% --- 摘要正文：小四宋体 ---
{
\songti\zihao{-4}
【在此处填写摘要正文，300-500字。】
\begin{enumerate}
    \item 问题背景与意义：一句话说明
    \item 主要方法：简述采用的模型和算法
    \item 关键结果：给出主要数值结果
    \item 创新点：本文贡献
\end{enumerate}

\vspace{0.8cm}
\noindent{\heiti\zihao{-4}关键词：}【关键词1】\quad 【关键词2】\quad 【关键词3】\quad 【关键词4】
}

\newpage

% ==========================================
% 正文开始
% ==========================================
\setcounter{page}{1}
\pagestyle{fancy}

% ==========================================
% 1. 问题重述
% ==========================================
\section{问题重述}

\subsection{问题背景}
【描述问题的实际背景和意义，约300-400字。】

\subsection{问题描述}
【用学术语言重述题目给出的问题。逐问说明已知条件、求解目标和约束。约400-500字。不要照抄原文。】

% ==========================================
% 2. 问题分析
% ==========================================
\section{问题分析}

\subsection{问题1的分析}
【分析核心难点、关键约束、解题思路。约200-300字。】

\subsection{问题2的分析}
【分析问题2与问题1的差异。约200-300字。】

\subsection{问题3的分析}
【分析动态特点。约200-300字。】

% ==========================================
% 3. 模型假设与符号说明
% ==========================================
\section{模型假设与符号说明}

\subsection{模型假设}
\begin{enumerate}[label=\arabic*.]
    \item 【假设1】
    \item 【假设2】
    \item 【假设3】
    \item 【假设4】
    \item 【假设5】
\end{enumerate}

\subsection{符号说明}
% 参数表
\begin{table}[H]
\centering
\caption{模型参数}
\begin{tabular}{c|l|c}
\toprule
\textbf{符号} & \textbf{含义} & \textbf{单位} \\
\midrule
$N$ & 【含义】 & 【单位】 \\
\bottomrule
\end{tabular}
\end{table}

% 决策变量表
\begin{table}[H]
\centering
\caption{决策变量}
\begin{tabular}{c|l|c}
\toprule
\textbf{符号} & \textbf{含义} & \textbf{类型} \\
\midrule
$x_{ij}$ & 【含义】 & 0-1变量 \\
\bottomrule
\end{tabular}
\end{table}

% ==========================================
% 4. 模型的建立与求解
% ==========================================
\section{模型的建立与求解}

\subsection{数据预处理}
\subsubsection{数据概况}
【描述数据来源、规模、特征。约300-400字。】

\subsubsection{数据处理}
【缺失值处理、异常值处理、特征工程等。约200-300字。】

\subsection{问题1的模型与求解}

\subsubsection{模型建立}

\textbf{目标函数：}

考虑【相关因素】，以【目标描述】最小为目标，建立如下模型：

\begin{equation}\label{eq:obj}
\min Z = \sum_{i} \sum_{j} c_{ij} x_{ij}
\end{equation}

其中，$c_{ij}$表示【含义】，$x_{ij}$表示【含义】。

\textbf{约束条件：}

\begin{enumerate}[label=(\arabic*)]
    \item \textbf{【约束名称】}：
    \begin{equation}\label{eq:con1}
        \sum_{j} x_{ij} \leq a_i, \quad \forall i
    \end{equation}
    约束(\ref{eq:con1})确保【实际含义】。如果不满足此约束，【说明后果】。

    \item \textbf{【约束名称】}：
    \begin{equation}\label{eq:con2}
        \sum_{i} x_{ij} \geq b_j, \quad \forall j
    \end{equation}
    约束(\ref{eq:con2})确保【实际含义】。
\end{enumerate}

\subsubsection{求解算法}

【算法选择理由】。

\begin{algorithm}[H]
\caption{【算法名称】}
\label{alg:1}
\begin{algorithmic}[1]
\REQUIRE 【输入参数】
\ENSURE 【输出结果】
\STATE 初始化
\FOR{【迭代条件】}
    \STATE 【步骤】
\ENDFOR
\RETURN 【结果】
\end{algorithmic}
\end{algorithm}

\textbf{算法关键设计：}
\begin{enumerate}
    \item \textbf{初始解生成}：【描述】
    \item \textbf{邻域结构}：【描述】
    \item \textbf{参数设置}：【描述】
\end{enumerate}

\subsubsection{求解结果}

【用表格展示结果，配合图表解读分析。】

\begin{table}[H]
\centering
\caption{问题1求解结果汇总}
\label{tab:q1-results}
\begin{tabular}{l|c|c}
\toprule
\textbf{指标} & \textbf{数值} & \textbf{单位} \\
\midrule
目标函数值 & 【】 & 【】 \\
\bottomrule
\end{tabular}
\end{table}

\subsection{问题2的模型与求解}
【结构同问题1，约800-1000字。】

\subsection{问题3的模型与求解}
【结构同问题1，约800-1000字。】

% ==========================================
% 5. 结果分析与检验
% ==========================================
\section{结果分析与检验}

\subsection{结果分析}
【逐项分析各问题的求解结果，配合图表。】

\subsection{灵敏度分析}
【关键参数的灵敏度分析。约400-500字。】

\subsection{模型检验}
【验证模型正确性和结果可信度的方法与结论。】

% ==========================================
% 6. 模型评价与推广
% ==========================================
\section{模型评价与推广}

\subsection{模型优点}
\begin{enumerate}
    \item 【优点1】
    \item 【优点2】
    \item 【优点3】
    \item 【优点4】
\end{enumerate}

\subsection{模型缺点}
\begin{enumerate}
    \item 【缺点1】
    \item 【缺点2】
    \item 【缺点3】
\end{enumerate}

\subsection{模型推广}
【模型在其他场景的应用前景。约300字。】

% ==========================================
% 参考文献
% ==========================================
\newpage
\begin{thebibliography}{99}
\bibitem{ref1} 作者1, 作者2. 文章标题[J]. 期刊名, 年份, 卷(期): 页码.
\bibitem{ref2} 作者. 文章标题[J]. 期刊名, 年份, 卷(期): 页码.
\bibitem{ref3} 作者. 书名[M]. 出版地: 出版社, 年份.
\end{thebibliography}

% ==========================================
% 附录
% ==========================================
\newpage
\section*{附录}
\addcontentsline{toc}{section}{附录}

\subsection*{附录1：核心求解代码}
{\small\ttfamily
\begin{verbatim}
# 在此粘贴核心代码
\end{verbatim}
}

\subsection*{附录2：补充数据}
【补充数据表格】

\end{document}
